% =====================================================================
%  CHAPTER 17: 学校健康教育（对应大纲第十章）
%  参考PPT：第十七章 学校健康教育（罗江洪，赣南医科大学）
%  往届笔记：第十五章 学校健康教育与健康促进
% =====================================================================
\chapter{学校健康教育}
\setcounter{chapter}{17}
\chapterintro{
\textbf{【掌握/双横线】}学校健康教育的概念；学校健康教育的基本内容（5个领域、5个水平）。

\textbf{【熟悉/单横线】}学校专题教育（预防艾滋病、毒品预防、预防烟草、学校性教育）；学校健康教育与健康促进的主要理论与应用（生活技能理论、组织改变理论、创新扩散理论、理性行动理论、大众意见领袖理论）；学校健康教育与健康促进的评价类型（形成评价、过程评价、效果评价）和评价方法（观察法、访谈法、调查问卷）。
}

% ----- 11.1 概述 -----
\section{概述}

\subsection{学校健康教育概念}

\begin{keybox}
\textbf{学校健康教育（School Health Education）：}以促进学生健康为核心的教育活动与过程。通过有计划、有组织、多种形式的教育教学活动，使学生掌握卫生保健知识，增强学生自我保健意识，养成科学、文明、健康的生活方式和行为习惯，达到预防疾病、增进身心健康、全面提高学生健康水平的目的。
\end{keybox}

\begin{defbox}
\textbf{健康促进学校（Health Promoting Schools）：}学校内所有成员为维护和促进师生健康共同努力，制定促进师生健康的规章制度，提供完整、积极的经验和知识结构，包括设置正式和非正式健康课程，创造安全健康的学校环境，提供适宜的卫生服务。

\textbf{健康素养（Health Literacy）：}个人获取和正确理解健康信息与服务，并运用这些信息和服务作出正确判断，维护和促进自身健康的能力。

\textbf{学生健康素养：}学生获取和正确理解健康信息与服务，并运用这些信息和服务作出正确判断，维护和促进自身健康的能力。
\end{defbox}

\subsection{学校健康教育基本内容}

\begin{keybox}
\textbf{5个领域：}
\begin{enumerate}[leftmargin=*]
    \item 健康行为与生活方式
    \item 疾病预防
    \item 心理健康
    \item 生长发育与青春期保健
    \item 安全应急与避险
\end{enumerate}

\textbf{5个水平：}
\begin{enumerate}[leftmargin=*]
    \item 水平一（小学1--2年级）
    \item 水平二（小学3--4年级）
    \item 水平三（小学5--6年级）
    \item 水平四（初中7--9年级）
    \item 水平五（高中10--12年级）
\end{enumerate}
\end{keybox}

% ----- 11.2 专题教育 -----
\section{专题教育}

\begin{defbox}
\textbf{学校专题健康教育主要包括：}
\begin{enumerate}[leftmargin=*]
    \item \imp{预防艾滋病专题教育}
    \item \imp{毒品预防专题教育}
    \item \imp{预防烟草教育}
\end{enumerate}
\end{defbox}

\subsection{学校性教育}

\begin{defbox}
\textbf{学校性教育（School-Based Sexuality Education）：}是以学校为场所，以学生为主要教育对象，通过有计划、有组织、有系统的教育活动进行的性科学教育。旨在帮助儿童青少年获得科学的性知识，形成正确的性态度和价值观，发展健康的人际交往能力，学会自我保护，促进身心健康发展。
\end{defbox}

% ----- 11.3 学校健康教育与健康促进实施 -----
\section{学校健康教育与健康促进实施}

\subsection{教学方法}

\begin{keybox}
\textbf{（一）直接式}

1. \imp{传统式}：课堂讲授；讲座；示教。

2. \imp{参与式}：小组讨论和案例分析；头脑风暴；角色扮演、小品和游戏活动；辩论和演讲；同伴教育。

\textbf{（二）间接式}

大众媒体；视听手段；网络系统性学习。
\end{keybox}

\subsection{学校健康教育组织实施}

\begin{keybox}
\begin{enumerate}[leftmargin=*]
    \item \imp{课堂教学}
    \item \imp{教师培训}
    \item \imp{学校支持性环境}
\end{enumerate}
\end{keybox}

% ----- 11.4 主要理论与应用 -----
\section{健康教育与健康促进主要理论与学校应用}

\begin{defbox}
\textbf{（一）生活技能理论}

\textbf{生活技能（Life Skills）：}一个人的心理社会能力。

\textbf{生活技能教育（Life Skills-Based Education）：}促进心理社会能力在适宜的文化背景下实践并得到发展，促进个体和社会发展，预防可能出现的健康和社会问题，保障人权。

\textbf{生活技能核心能力（5对10项）：}
\begin{enumerate}[leftmargin=*]
    \item \imp{自我认识能力——同理能力}
    \item \imp{有效交流能力——人际关系能力}
    \item \imp{调节情绪能力——缓解压力能力}
    \item \imp{创造性思维能力——批判性思维能力}
    \item \imp{决策能力——解决问题能力}
\end{enumerate}
\end{defbox}

\begin{defbox}
\textbf{（二）组织改变理论}

\textbf{组织改变理论模型（Theory of Organizational Change）：}通过人群所在组织的改变、规章制度和管理模式的完善，实现对组织内全体人群的干预。

\textbf{4个阶段：}
\begin{enumerate}[leftmargin=*]
    \item \imp{确立问题或认知阶段}：对问题的认知与分析，寻求解决问题的方法和评价
    \item \imp{初期行动阶段}：制订项目规划和实施计划，为开始改变准备
    \item \imp{执行阶段}：创新的执行阶段
    \item \imp{制度化阶段}：新政策等变为确定的、新的目标价值融入组织里
\end{enumerate}
\end{defbox}

\begin{defbox}
\textbf{（三）创新扩散理论}

\textbf{创新扩散（Diffusion of Innovation, DI）：}一项新事物通过一定传播渠道在整个社区或某个人群内扩散，逐渐为社区成员或该人群成员了解与采用的过程。

\textbf{5个阶段：}
\begin{enumerate}[leftmargin=*]
    \item \imp{认知阶段}：个体首次接触新事物，但是缺少相关信息
    \item \imp{说服阶段}：个体对新事物产生兴趣，积极寻求相关信息
    \item \imp{决策阶段}：评估使用新事物的优势和劣势，决定采用或拒绝
    \item \imp{实施阶段}：个体不同程度采用了新事物
    \item \imp{确认阶段}：个人定型，决定继续使用创新
\end{enumerate}

\textbf{5种类型：}
\begin{enumerate}[leftmargin=*]
    \item \imp{创新者}：人群中最先接受信息
    \item \imp{早期采用者}：较易接受新观念、态度慎重、常具领导力
    \item \imp{中期采用者}：慎重、深思熟虑
    \item \imp{后期采用者}：倾向于对创新事物持怀疑态度，等多数人接受并认同该创新才采用
    \item \imp{迟缓者}：观念较保守、坚持习惯，不到万不得已不愿接受创新
\end{enumerate}
\end{defbox}

\begin{defbox}
\textbf{（四）理性行动理论}

\textbf{理性行动理论（Theory of Reasoned Action, TRA）/ 理性行为理论：}分析态度如何有意识影响个体行为的理论。
\end{defbox}

\begin{defbox}
\textbf{（五）大众意见领袖理论}

\textbf{大众意见领袖（Popular Opinion Leaders, POLs）/ 舆论领袖：}大众信息传播中信息从源头到一般受众的中间处理环节或节点。

\textbf{信息传播模式：}大众传播——意见领袖——一般受众。
\end{defbox}

% ----- 11.5 评价 -----
\section{学校健康教育与健康促进评价}

\subsection{评价类型}

\begin{defbox}
\textbf{（一）形成评价}

\textbf{（二）过程评价}
\begin{enumerate}[leftmargin=*]
    \item 健康教育活动评价
    \item 学校卫生服务评价
    \item 学校环境评价
\end{enumerate}

\textbf{（三）效果评价}
\begin{enumerate}[leftmargin=*]
    \item 健康教育活动效果评价
    \item 卫生服务效果评价
\end{enumerate}
\end{defbox}

\subsection{评价方法}

\begin{defbox}
\begin{enumerate}[leftmargin=*]
    \item \imp{观察法}
    \item \imp{个人访谈、小组访谈和资料与案例分析}
    \item \imp{调查问卷}：设置封闭式问题和开放式问题，可用于知识、态度和行为的评价
\end{enumerate}
\end{defbox}

\subsection{评价维度与指标}

\begin{defbox}
学校健康教育与健康促进评价主要围绕四个维度展开：

\textbf{知识、态度、技能、行为。}
\end{defbox}

\subsection{RE-AIM评价框架}

\begin{keybox}
\textbf{RE-AIM框架}是学校健康教育项目评价的综合框架，由五个维度组成：

\begin{enumerate}[leftmargin=*]
    \item \imp{可及性（Reach）}：项目覆盖目标人群的比例及代表性。评价指标包括参与率、覆盖范围、参与者与未参与者的特征比较。

    \item \imp{有效性（Effectiveness）}：项目在关键健康指标上的实际效果。包括知识、态度、技能、行为等维度的改变，以及正性和负性结果的全面评估。

    \item \imp{采纳性（Adoption）}：实施项目的机构和人员采纳干预措施的程度。评价学校、教师等层面对项目的接受和参与情况。

    \item \imp{实施性（Implementation）}：项目按原计划执行的程度和保真度。包括各组成部分的实施质量、时间投入、成本等过程指标。

    \item \imp{维持性（Maintenance）}：项目效果在个体和组织层面的长期持续性。个体层面指行为改变的长期保持，组织层面指项目在机构中的制度化程度。
\end{enumerate}

\textbf{RE-AIM框架的意义：}通过五个维度的综合评价，不仅可以了解项目"是否有效"，还能评估其"在多大范围内有效""在什么条件下有效""能否持续推广"，为学校健康教育项目的设计、实施和推广提供全面指导。该框架弥补了传统评价仅关注有效性的不足，是当前国际上学校健康促进项目评价的推荐框架。
\end{keybox}

\newpage
